% ============================================================================
%  Conclusion -- % ============================================================================
%  Conclusion -- % ============================================================================
%  Conclusion -- % ============================================================================
%  Conclusion -- \input{sections/conclusion}
%  Sentences tied to results use \res* macros; the \todo notes mark findings to state
%  (or delete) once Tables in sections/results.tex are filled.
% ============================================================================
\section{Conclusion}
\label{sec:conclusion}

We presented HyperAVNet, a hypercomplex audio-visual network for fish feeding intensity assessment.
Hamilton-product weight sharing runs through every parameterised part of the model: the backbones,
attention, fusion, feedback and heads. Fusion is expressed as an algebraic operation: hypercomplex
component fusion treats the audio view, the video view and their agreement as the components of one
quaternion or octonion, and mixes them with a single shared set of weights, at a parameter cost that
does not grow with the number of interaction terms. Because the fusion blocks can be placed at any
depth, with zero-initialised feedback into both streams, the question of \emph{where} to fuse becomes a
controlled experiment rather than a design assumption.

On AV-FFIA, the distilled quaternion model reaches \resOurAcc\% accuracy (macro-F1 \resOurFone\%) with
\resOurParams\,M parameters and \resOurGMACs\,GMACs per two-second clip. That is an order of
magnitude fewer parameters and operations than the U-FFIA models. \todo{State the position relative to
DynaGate-AV and HMMF-Net under their protocols.} The ablations show that
\todo{(i) the best fusion set was $\mathcal{S}=\resBestDepth$ (late fusion: \resLateAcc\%; with
feedback: \resProgAcc\%; without feedback: \resNoFeedbackAcc\%);}
\todo{(ii) the quaternion and octonion networks keep the accuracy of the real network (\resRealAcc\%)
with $3.7\times$ and $6.4\times$ fewer parameters (\resQuatAcc\% and \resOctAcc\%), or say where they
fall short;}
\todo{(iii) HCF compares with the attention bottleneck as \resHCFAcc\% vs.\ \resMBTAcc\%;}
and \todo{(iv) distillation changes the student's accuracy by \resKDGain\ points, with the teacher
ablation showing whether teacher strength or teacher type matters more.}
On the session-disjoint split, accuracy is \resGroupedAcc\%, against \resOfficialAcc\% on the official
random split. \todo{If the gap is substantial, conclude that clip-level splits overestimate
generalisation on AV-FFIA, and recommend session-disjoint evaluation as the default for future work.}

\paragraph{Limitations and future work.}
The hypercomplex structure reduces stored parameters, memory and transmission size, but not
multiply--accumulate operations; further compute savings would require structured sparsity or
algebra-aware kernels. AV-FFIA covers one species in one controlled tank. Validation on farm data, such
as the turbid, foamy tilapia tanks in which U-FFIA was field-tested~\cite{cui2024uffia}, and on other
species is the most important next step. Three further directions follow naturally from this design:
a streaming, causal variant for continuous monitoring and closed-loop feeding control; hypercomplex
teachers and self-supervised audio-visual pretraining on unlabelled farm recordings; and extending HCF
to additional sensors (e.g.\ sonar or water-quality signals), which only requires a larger algebra
(octonion or sedenion components) rather than a wider fusion layer. To make comparisons on AV-FFIA
reproducible, we release our code, the split files for all three protocols with their hashes, and the
trained models.

%  Sentences tied to results use \res* macros; the \todo notes mark findings to state
%  (or delete) once Tables in sections/results.tex are filled.
% ============================================================================
\section{Conclusion}
\label{sec:conclusion}

We presented HyperAVNet, a hypercomplex audio-visual network for fish feeding intensity assessment.
Hamilton-product weight sharing runs through every parameterised part of the model: the backbones,
attention, fusion, feedback and heads. Fusion is expressed as an algebraic operation: hypercomplex
component fusion treats the audio view, the video view and their agreement as the components of one
quaternion or octonion, and mixes them with a single shared set of weights, at a parameter cost that
does not grow with the number of interaction terms. Because the fusion blocks can be placed at any
depth, with zero-initialised feedback into both streams, the question of \emph{where} to fuse becomes a
controlled experiment rather than a design assumption.

On AV-FFIA, the distilled quaternion model reaches \resOurAcc\% accuracy (macro-F1 \resOurFone\%) with
\resOurParams\,M parameters and \resOurGMACs\,GMACs per two-second clip. That is an order of
magnitude fewer parameters and operations than the U-FFIA models. \todo{State the position relative to
DynaGate-AV and HMMF-Net under their protocols.} The ablations show that
\todo{(i) the best fusion set was $\mathcal{S}=\resBestDepth$ (late fusion: \resLateAcc\%; with
feedback: \resProgAcc\%; without feedback: \resNoFeedbackAcc\%);}
\todo{(ii) the quaternion and octonion networks keep the accuracy of the real network (\resRealAcc\%)
with $3.7\times$ and $6.4\times$ fewer parameters (\resQuatAcc\% and \resOctAcc\%), or say where they
fall short;}
\todo{(iii) HCF compares with the attention bottleneck as \resHCFAcc\% vs.\ \resMBTAcc\%;}
and \todo{(iv) distillation changes the student's accuracy by \resKDGain\ points, with the teacher
ablation showing whether teacher strength or teacher type matters more.}
On the session-disjoint split, accuracy is \resGroupedAcc\%, against \resOfficialAcc\% on the official
random split. \todo{If the gap is substantial, conclude that clip-level splits overestimate
generalisation on AV-FFIA, and recommend session-disjoint evaluation as the default for future work.}

\paragraph{Limitations and future work.}
The hypercomplex structure reduces stored parameters, memory and transmission size, but not
multiply--accumulate operations; further compute savings would require structured sparsity or
algebra-aware kernels. AV-FFIA covers one species in one controlled tank. Validation on farm data, such
as the turbid, foamy tilapia tanks in which U-FFIA was field-tested~\cite{cui2024uffia}, and on other
species is the most important next step. Three further directions follow naturally from this design:
a streaming, causal variant for continuous monitoring and closed-loop feeding control; hypercomplex
teachers and self-supervised audio-visual pretraining on unlabelled farm recordings; and extending HCF
to additional sensors (e.g.\ sonar or water-quality signals), which only requires a larger algebra
(octonion or sedenion components) rather than a wider fusion layer. To make comparisons on AV-FFIA
reproducible, we release our code, the split files for all three protocols with their hashes, and the
trained models.

%  Sentences tied to results use \res* macros; the \todo notes mark findings to state
%  (or delete) once Tables in sections/results.tex are filled.
% ============================================================================
\section{Conclusion}
\label{sec:conclusion}

We presented HyperAVNet, a hypercomplex audio-visual network for fish feeding intensity assessment.
Hamilton-product weight sharing runs through every parameterised part of the model: the backbones,
attention, fusion, feedback and heads. Fusion is expressed as an algebraic operation: hypercomplex
component fusion treats the audio view, the video view and their agreement as the components of one
quaternion or octonion, and mixes them with a single shared set of weights, at a parameter cost that
does not grow with the number of interaction terms. Because the fusion blocks can be placed at any
depth, with zero-initialised feedback into both streams, the question of \emph{where} to fuse becomes a
controlled experiment rather than a design assumption.

On AV-FFIA, the distilled quaternion model reaches \resOurAcc\% accuracy (macro-F1 \resOurFone\%) with
\resOurParams\,M parameters and \resOurGMACs\,GMACs per two-second clip. That is an order of
magnitude fewer parameters and operations than the U-FFIA models. \todo{State the position relative to
DynaGate-AV and HMMF-Net under their protocols.} The ablations show that
\todo{(i) the best fusion set was $\mathcal{S}=\resBestDepth$ (late fusion: \resLateAcc\%; with
feedback: \resProgAcc\%; without feedback: \resNoFeedbackAcc\%);}
\todo{(ii) the quaternion and octonion networks keep the accuracy of the real network (\resRealAcc\%)
with $3.7\times$ and $6.4\times$ fewer parameters (\resQuatAcc\% and \resOctAcc\%), or say where they
fall short;}
\todo{(iii) HCF compares with the attention bottleneck as \resHCFAcc\% vs.\ \resMBTAcc\%;}
and \todo{(iv) distillation changes the student's accuracy by \resKDGain\ points, with the teacher
ablation showing whether teacher strength or teacher type matters more.}
On the session-disjoint split, accuracy is \resGroupedAcc\%, against \resOfficialAcc\% on the official
random split. \todo{If the gap is substantial, conclude that clip-level splits overestimate
generalisation on AV-FFIA, and recommend session-disjoint evaluation as the default for future work.}

\paragraph{Limitations and future work.}
The hypercomplex structure reduces stored parameters, memory and transmission size, but not
multiply--accumulate operations; further compute savings would require structured sparsity or
algebra-aware kernels. AV-FFIA covers one species in one controlled tank. Validation on farm data, such
as the turbid, foamy tilapia tanks in which U-FFIA was field-tested~\cite{cui2024uffia}, and on other
species is the most important next step. Three further directions follow naturally from this design:
a streaming, causal variant for continuous monitoring and closed-loop feeding control; hypercomplex
teachers and self-supervised audio-visual pretraining on unlabelled farm recordings; and extending HCF
to additional sensors (e.g.\ sonar or water-quality signals), which only requires a larger algebra
(octonion or sedenion components) rather than a wider fusion layer. To make comparisons on AV-FFIA
reproducible, we release our code, the split files for all three protocols with their hashes, and the
trained models.

%  Sentences tied to results use \res* macros; the \todo notes mark findings to state
%  (or delete) once Tables in sections/results.tex are filled.
% ============================================================================
\section{Conclusion}
\label{sec:conclusion}

We presented HyperAVNet, a hypercomplex audio-visual network for fish feeding intensity assessment.
Hamilton-product weight sharing runs through every parameterised part of the model: the backbones,
attention, fusion, feedback and heads. Fusion is expressed as an algebraic operation: hypercomplex
component fusion treats the audio view, the video view and their agreement as the components of one
quaternion or octonion, and mixes them with a single shared set of weights, at a parameter cost that
does not grow with the number of interaction terms. Because the fusion blocks can be placed at any
depth, with zero-initialised feedback into both streams, the question of \emph{where} to fuse becomes a
controlled experiment rather than a design assumption.

On AV-FFIA, the distilled quaternion model reaches \resOurAcc\% accuracy (macro-F1 \resOurFone\%) with
\resOurParams\,M parameters and \resOurGMACs\,GMACs per two-second clip. That is an order of
magnitude fewer parameters and operations than the U-FFIA models. \todo{State the position relative to
DynaGate-AV and HMMF-Net under their protocols.} The ablations show that
\todo{(i) the best fusion set was $\mathcal{S}=\resBestDepth$ (late fusion: \resLateAcc\%; with
feedback: \resProgAcc\%; without feedback: \resNoFeedbackAcc\%);}
\todo{(ii) the quaternion and octonion networks keep the accuracy of the real network (\resRealAcc\%)
with $3.7\times$ and $6.4\times$ fewer parameters (\resQuatAcc\% and \resOctAcc\%), or say where they
fall short;}
\todo{(iii) HCF compares with the attention bottleneck as \resHCFAcc\% vs.\ \resMBTAcc\%;}
and \todo{(iv) distillation changes the student's accuracy by \resKDGain\ points, with the teacher
ablation showing whether teacher strength or teacher type matters more.}
On the session-disjoint split, accuracy is \resGroupedAcc\%, against \resOfficialAcc\% on the official
random split. \todo{If the gap is substantial, conclude that clip-level splits overestimate
generalisation on AV-FFIA, and recommend session-disjoint evaluation as the default for future work.}

\paragraph{Limitations and future work.}
The hypercomplex structure reduces stored parameters, memory and transmission size, but not
multiply--accumulate operations; further compute savings would require structured sparsity or
algebra-aware kernels. AV-FFIA covers one species in one controlled tank. Validation on farm data, such
as the turbid, foamy tilapia tanks in which U-FFIA was field-tested~\cite{cui2024uffia}, and on other
species is the most important next step. Three further directions follow naturally from this design:
a streaming, causal variant for continuous monitoring and closed-loop feeding control; hypercomplex
teachers and self-supervised audio-visual pretraining on unlabelled farm recordings; and extending HCF
to additional sensors (e.g.\ sonar or water-quality signals), which only requires a larger algebra
(octonion or sedenion components) rather than a wider fusion layer. To make comparisons on AV-FFIA
reproducible, we release our code, the split files for all three protocols with their hashes, and the
trained models.
